% Separate post-CP18 research. Candidate; all-family hypotheses are explicit.
\subsection{Artinian realization with a retained full second neighbourhood}

\begin{theorem}[Candidate: cohomologically flat full-neighbourhood realization]
\label{D:Artinian-full-neighbourhood-realization}
Let $A$ be a complete excellent normal local surface with
algebraically closed coefficient and residue field $k$. Let
$\pi:Y\to\operatorname{Spec}A$ be a projective regular resolution,
an isomorphism on the puncture, and let $Z$ be a reduction cycle
of full exceptional support in the sense of
Theorem~\ref{D:reduction-cycle-full-module-correspondence}.
The characteristic of $k$ is arbitrary.

Let $T$ be a finite-dimensional local Artinian $k$-algebra with
residue field $k$, put $A_T=A\otimes_kT$ and $Y_T=Y\times_kT$,
and let $V_2$ be a vector bundle on $(2Z)_T$. Write
$V=V_2|_{Z_T}$ and $V_0=V\otimes_Tk$. Assume:
\begin{enumerate}
\item $V_0$ is generically generated by its global sections;
\item for $W=V$ and $W=V(Z)$, both $H^0(Z_T,W)$ and
$H^1(Z_T,W)$ are finite free over $T$, and their formation
commutes with every homomorphism from $T$ to a finite local
Artinian $k$-algebra with residue field $k$;
\item the actual sequence
\[
 0\longrightarrow V\longrightarrow V_2(Z)
       \longrightarrow V(Z)\longrightarrow0
\]
has a split injective boundary
\[
 \delta_0:H^0(Z_T,V(Z))\longrightarrow H^1(Z_T,V).
                                                        \tag{D.AR.1}
\]
\end{enumerate}
There is a vector bundle $\mathcal F_T$ on $Y_T$ restricting
to the given $V_2$. Every such vector-bundle extension has the
following properties. Its module
$M_T=H^0(Y_T,\mathcal F_T)$ is finite over $A_T$, is flat
over $T$, and has maximal Cohen--Macaulay closed fibre.
There are natural identifications
\[
 \begin{split}
 M_T&=H^0(Y_T\setminus Z_T,\mathcal F_T),\\
 (\pi_T^*M_T)^{\vee\vee}&\simeq\mathcal F_T ,
 \end{split}                                           \tag{D.AR.2}
\]
where the double dual is taken over $\mathcal O_{Y_T}$.
Moreover, for every homomorphism $T\to T'$ as in item~(2),
\[
 M_T\otimes_TT'\ \simeq\
 H^0(Y_{T'},\mathcal F_T\otimes_TT').                    \tag{D.AR.3}
\]
In particular these are actual relative maximal
Cohen--Macaulay modules with their full vector-bundle pullbacks,
not only fibrewise realizations.

The hypotheses retain the second neighbourhood and the
cohomology base-change conditions. The theorem does not
assert that arbitrary families of bundles on $Z$ satisfy them,
that a prescribed map on $Z$ lifts, or that a choice of
extension $\mathcal F_T$ is functorial in $V_2$.
\end{theorem}

\begin{proof}
All thickenings and coefficient base changes in this proof
are over the fixed $k$-scheme $Y$. The ring $A_T$ is complete
Noetherian and local: as an $A$-module it is a finite direct
sum of copies of $A$, and its nilpotent coefficient ideal has
quotient $A$. The ideals of $Z_T$ and of the closed fibre of
$Y_T\to\operatorname{Spec}A_T$ have the same radical, so their
adic topologies are cofinal.

\emph{Existence of a vector-bundle extension.}
If $V_j$ is already constructed on $(jZ)_T$, the obstruction
to extending it across $(jZ)_T\subset((j+1)Z)_T$ lies in
\[
 H^2(Z_T,\mathcal End(V)(-jZ))=0.
                                                        \tag{D.AR.4}
\]
Indeed local free lifts exist, and the discrepancy of their
transition matrices is the usual additive degree-two cocycle
in the displayed square-zero ideal. The vanishing holds
because $Z_T$ has dimension one. This is also the
$\operatorname{Ext}^2$ obstruction of
\citep[Lemma~91.10.4 and complete proof, Tag 08MF]
{StacksRelativeKahnB}.
Starting from the prescribed $V_2$, choose compatible lifts
for every $j$. Projective Grothendieck existence algebraizes
the resulting formal coherent sheaf
\citep[Lemma~30.24.3 and complete proof, Tag 0885]
{StacksRelativeKahnB}. At each point of the closed fibre its
completion is free, so the algebraized sheaf is locally free
there. Its nonfree locus is closed and proper over the local
base and misses the closed fibre. It is therefore empty.
This gives $\mathcal F_T$.

\emph{Cohomology away from the first neighbourhood.}
The generic-generation lemma in the proof of
Theorem~\ref{D:reduction-cycle-full-module-correspondence}
gives
\[
                  H^1(Y,\mathcal F_0(-Z))=0,
                                                        \tag{D.AR.5}
\]
where $\mathcal F_0=\mathcal F_T\otimes_Tk$.
For clarity, it only uses generic generation of $V_0$,
generic generation of $\mathcal O_Z(-Z)$, and
$H^1(Z,\mathcal O_Z(-Z))=0$; fullness is not assumed at this
step. An Artinian filtration of any coefficient algebra $T'$
has successive quotients copies of $k$. Since
$\mathcal F_T\otimes_TT'$ is a vector bundle, (D.AR.5) and
the associated exact sequences give
\[
 H^1(Y_{T'},(\mathcal F_T\otimes_TT')(-Z))=0.
                                                        \tag{D.AR.6}
\]
All higher cohomology of coherent sheaves on $Y_{T'}$
vanishes in degrees at least two. One may see this by formal
functions: the closed-fibre thickenings have dimension one,
and the finite higher-cohomology modules over the complete
local base have zero completion. Hence restriction gives
\[
 H^1(Y_{T'},\mathcal F_T\otimes_TT')
             \simeq H^1(Z_{T'},V\otimes_TT').
                                                        \tag{D.AR.7}
\]

\emph{Flatness and base change of the section module.}
A finite affine cover of the separated scheme $Y_T$
computes the cohomology of $\mathcal F_T(-Z)$ by a bounded
Čech complex of flat $T$-modules. Equation (D.AR.6) says that
this complex and its coefficient base changes have no
positive cohomology. Equivalently, its degree-zero
cohomology is flat and commutes with these base changes.
This also follows directly by descending through the bounded
Čech complex: successive kernels of surjections between
flat modules are flat. The restriction sequence and
(D.AR.6) now give
\[
 0\longrightarrow H^0(Y_T,\mathcal F_T(-Z))
 \longrightarrow M_T
 \longrightarrow H^0(Z_T,V)\longrightarrow0.
                                                        \tag{D.AR.8}
\]
The right term is free and commutes with coefficient change
by assumption. Both conclusions therefore hold for $M_T$
as well. This proves flatness and (D.AR.3). Finiteness over
$A_T$ follows from properness.

\emph{The full Hartogs property over $T$.}
The split injection (D.AR.1) remains injective after every
coefficient change under consideration. For each $j\ge1$,
the reduction-cycle generic-generation lemma supplies a
fixed map on $Y$,
\[
 \mathcal O_Y^a\longrightarrow\mathcal O_Y(-jZ),
\]
which is onto at the exceptional generic points. Its dual
restricts to an injection
$\mathcal O_Z(jZ)\to\mathcal O_Z^a$. This injection remains
injective after coefficient extension, since that extension
is flat over $k$. Multiplication gives a commuting comparison
of the boundary
\[
 \delta_j:H^0(Z_T,V((j+1)Z))
                \longrightarrow H^1(Z_T,V(jZ))
                                                        \tag{D.AR.9}
\]
with $a$ copies of (D.AR.1). The comparison is defined on
$(2Z)_T$, and its map on the degree-zero source is injective.
Thus every $\delta_j$ is injective.

Put $F_j=\mathcal F_T|_{(jZ)_T}$. The sequences
\[
 0\longrightarrow F_j(jZ)
 \longrightarrow F_{j+1}((j+1)Z)
 \longrightarrow V((j+1)Z)\longrightarrow0
                                                        \tag{D.AR.10}
\]
show that the transition
$H^0(F_j(jZ))\to H^0(F_{j+1}((j+1)Z))$
is an isomorphism: the next boundary, composed with
restriction, is the injective map (D.AR.9).
Consequently
\[
 H^1_{Z_T}(Y_T,\mathcal F_T)
  =\varinjlim_j H^0((jZ)_T,F_j(jZ))
  =H^0(Z_T,V(Z)).
                                                        \tag{D.AR.11}
\]
The first equality is the Cartier local-cohomology
calculation; $\mathcal F_T$ has no $Z_T$-supported torsion.
Under (D.AR.7), the map from (D.AR.11) to
$H^1(Y_T,\mathcal F_T)$ is exactly $\delta_0$.
It is injective. The localization exact sequence now gives
the first equality in (D.AR.2).

On the closed fibre, the same argument and generic generation
give the ordinary full-bundle criterion. Thus
$M_0=H^0(Y,\mathcal F_0)$ is maximal Cohen--Macaulay and
$\mathcal F_0=(\pi^*M_0)^{\vee\vee}$.
There is no hidden flatness assertion here:
$M_T\otimes_Tk=M_0$ was already proved in (D.AR.3).
For any finite regular Noether normalization
$S=k[[u,v]]\subset A$, the $S$-module $M_0$ is free.
Lift an $S$-basis to $M_T$. Nilpotent Nakayama gives a
surjection $S_T^d\to M_T$; flatness over $T$ makes its kernel
reduce to zero, and nilpotent Nakayama kills the kernel.
Thus $M_T$ is $S_T$-free. In particular $u,v$ form a regular
sequence on it, as required for relative Hartogs.

\emph{Identification of the relative full pullback.}
Assumption~(2) and (D.AR.8) lift a finite set of generic
generators of $V_0$ to global sections of $\mathcal F_T$.
Nilpotent Nakayama implies that
$\pi_T^*M_T\to\mathcal F_T$ is surjective at every exceptional
generic point. Away from the exceptional locus it is an
isomorphism. Let $G$ be its image. The cokernel
$\mathcal F_T/G$ is supported in codimension two, and the
kernel of $\pi_T^*M_T\to G$ is killed locally by a power of
the fixed Cartier equation of $Z$.
The scheme $Y_T$ is Cohen--Macaulay of dimension two:
its local rings are finite free over the corresponding
regular local rings of $Y$. Hence a codimension-two
coherent module has zero $\operatorname{Ext}^0$ and
$\operatorname{Ext}^1$ into $\mathcal O_{Y_T}$.
The exceptional torsion kernel has no homomorphism into
$\mathcal O_{Y_T}$, since the Cartier equation is a
nonzerodivisor. Dualizing the two kernel/image sequences
therefore identifies
$(\pi_T^*M_T)^\vee$ with $\mathcal F_T^\vee$.
Dualizing again proves the second equality in (D.AR.2).
The proof applies to every chosen extension and every stated
coefficient base change. It never invokes a prescribed
$H^0$-isomorphism lift from Kahn's tangent-space comparison.
\end{proof}

\begin{corollary}[Candidate: exact two-level recovery for the realized Artinian families]
\label{D:Artinian-canonical-two-level-jets}
Use the fixed $R$, $J$, separating normalization
$S=k[[x,y]]$, scalar $c_0$, and integer $c$ of
Lemma~\ref{D:canonical-pullback-uniform-jet-loss}.
Let $A/R$ be a finite normal cover domain, étale on the
puncture, and let $Y$ be a projective regular resolution.
Let $(M_T,\mathcal F_T)$ be one of the realized families
constructed in
Theorem~\ref{D:Artinian-full-neighbourhood-realization}
on this resolution. In particular its evaluation map is
surjective at every exceptional generic point, its closed
fibre is the canonical full pullback of $M_0$, and
(D.AR.2)--(D.AR.3) hold.

Set
\[
 I_b=(x^b,y^b)A_T,\qquad
 Q_{b,T}=V(I_b\mathcal O_{Y_T}),\qquad
 V_b(T)=H^0(Q_{b,T},\mathcal F_T|_{Q_{b,T}}),
\]
and let $W_b(T)$ be the image of $V_{b+1}(T)\to V_b(T)$.
For $b\ge2\ell-1+2c$ there is a canonical isomorphism
\[
             W_b(T)/I_\ell W_b(T)\ \simeq\
                         M_T/I_\ell M_T.              \tag{D.AR.12}
\]
It is functorial for actual module maps and commutes with
every Artinian coefficient change in
Theorem~\ref{D:Artinian-full-neighbourhood-realization}.
The individual spaces $V_b(T)$ and the raw images $W_b(T)$
are not asserted to commute with base change.
\end{corollary}

\begin{proof}
We give the relative details rather than infer them from the
pointwise theorem. The module $M_T$ is free over $S_T$ by
the lifted-basis argument above, and relative Hartogs gives
$H^0(Y_T,\mathcal F_T)=M_T$.

The actual Noether-different factorizations of
Proposition~\ref{D:cover-uniform-Noether-different}
hold for these $T$-flat sources. In particular the fixed
$z\in\mathfrak N(R/S)$ used to define $c_0$ factors
$z\operatorname{id}_{M_T}$ through a finite free
$A_T$-module. This induces a factorization
\[
       \mathcal F_T\longrightarrow\mathcal O_{Y_T}^a
                          \longrightarrow\mathcal F_T
                                                        \tag{D.AR.13}
\]
with composite $z$. To check the first arrow, pull it back
from $M_T$ where evaluation is onto. The kernel of evaluation
is exceptional torsion, so its map to the free target
vanishes. Thus the first arrow is defined at every
codimension-one point. Hartogs for the locally free source
and target on the Cohen--Macaulay scheme $Y_T$ extends it
across the remaining codimension-two set. The composite is
$z$ since this is true on the puncture.

Put $K_b(T)=H^0(Y_T,I_b\mathcal O_{Y_T}\mathcal F_T)$.
The scalar clearing inclusion in the proof of
Lemma~\ref{D:canonical-pullback-uniform-jet-loss} is an
inclusion of fixed $k$-vector spaces. Flat coefficient
extension over $k$ gives
\[
 zH^0(Y_T,I_b\mathcal O_{Y_T})
                   \subset (x,y)^bA_T.
                                                        \tag{D.AR.14}
\]
Apply (D.AR.13) to a section of $K_b(T)$ and then apply
(D.AR.14) to its coordinates. Exactly as in the pointwise
argument,
$c_0^2K_b(T)\subset(x,y)^bM_T$.
Since $M_T$ is $S_T$-free and the initial form of the fixed
$c_0\in S$ remains a nonzerodivisor in $T[X,Y]$,
\[
 K_b(T)\subset(x,y)^{b-2c}M_T\subset I_\ell M_T
 \quad\text{if }b\ge2\ell-1+2c.                        \tag{D.AR.15}
\]

It remains to verify the adjacent-level cohomology
transition. Choose $r$ closed-fibre sections of $M_0$
that generate $\mathcal F_0$ at each exceptional generic
point and lift them to $M_T$. Their determinant is a
nonzerodivisor on $Y_T$: this follows by the Artinian
filtration from the same assertion on $Y$. Hence there is
an injection $\mathcal O_{Y_T}^r\to\mathcal F_T$.
Its cokernel has support finite over $\operatorname{Spec}A_T$,
because its underlying support misses the exceptional
generic points and is proper with finite fibres. Therefore
\[
 H^1(Y,\mathcal O_Y)^r\otimes_kT
                       \longrightarrow H^1(Y_T,\mathcal F_T)
\]
is surjective. The fixed canonical-defect ideal $J$ kills
the source, so $JH^1(Y_T,\mathcal F_T)=0$.

Let $D=\gcd(\operatorname{div}_Yx,\operatorname{div}_Yy)$.
The two-generator syzygy on $Y$, pulled back over $k$ and
then tensored with the vector bundle $\mathcal F_T$, is
\[
 0\longrightarrow\mathcal F_T(bD)
 \xrightarrow{(-y^b,x^b)}\mathcal F_T^2
 \xrightarrow{(x^b,y^b)}
 I_b\mathcal O_{Y_T}\mathcal F_T\longrightarrow0.
                                                        \tag{D.AR.16}
\]
The comparison from level $b+1$ to $b$ is multiplication
by $xy$ on the first term and $\operatorname{diag}(x,y)$
on the middle term. Since $x,y\in J$ and coherent $H^2$
vanishes on $Y_T$, it follows that
\[
 H^1(Y_T,I_{b+1}\mathcal O_{Y_T}\mathcal F_T)
   \longrightarrow H^1(Y_T,I_b\mathcal O_{Y_T}\mathcal F_T)
 \quad\text{is zero}.                                 \tag{D.AR.17}
\]
The connecting maps for restriction to $Q_{b,T}$ now give,
with no image/base-change interchange,
\[
          W_b(T)=\operatorname{im}(M_T\to V_b(T))
                         =M_T/K_b(T).
                                                        \tag{D.AR.18}
\]
Equations (D.AR.15) and (D.AR.18) prove (D.AR.12).
Functoriality follows from the actual evaluation and
restriction maps. The right side of (D.AR.12) commutes
with coefficient change, proving the stated base-change
assertion for the recovered quotient even when the raw
cohomology spaces and images do not.
\end{proof}

\begin{remark}[What is still required for the quantitative atlas]
The preceding results turn the specified cohomologically
flat $2Z$ data into actual Artinian families and recover
their canonical module jets. They do not provide a
polynomial presentation, with the desired degree-linear
coordinate budget, of the cohomology base-change strata,
their moving determinants, or all compatible Frobenius
arrows. In particular, an arbitrary map on $Z$ is not
declared liftable. No representation-growth or
finite-generation conclusion is inferred here.
\end{remark}
